\documentclass[conference]{IEEEtran}
\IEEEoverridecommandlockouts
\usepackage{cite}
\usepackage{amsmath,amssymb,amsfonts}
\usepackage{algorithmic}
\usepackage{algorithm}
\usepackage{graphicx}
\usepackage{textcomp}
\usepackage{xcolor}
\usepackage{booktabs}
\usepackage{hyperref}
\usepackage{url}

\begin{document}

\title{Unsupervised Neural Network for\\Multi-Genre Music Generation\\
{\large CSE425 Neural Networks --- Spring 2026}}

\author{\IEEEauthorblockN{Ummay Maimona Chaman}
\IEEEauthorblockA{\textit{ID: 22301719, Section: 1}\\
\textit{BRAC University}\\
Dhaka, Bangladesh\\
ummay.maimona.chaman@g.bracu.ac.bd}}

\maketitle

%------------------------------------------------------------------
\begin{abstract}
This paper presents a comprehensive study of unsupervised deep
learning architectures for automatic multi-genre music generation.
We design, implement, and evaluate four progressively advanced
systems: (1)~an LSTM Autoencoder for single-genre reconstruction,
(2)~a genre-conditioned Variational Autoencoder (VAE) for diverse
multi-genre synthesis, (3)~a decoder-only Transformer with causal
self-attention for long-horizon coherent sequence generation, and
(4)~a Reinforcement Learning from Human Feedback (RLHF) fine-tuned
generator. Experiments use the MAESTRO, Lakh MIDI, and Groove MIDI
datasets. Results demonstrate consistent improvement from Task~1 to
Task~4 across pitch histogram similarity, rhythm diversity, and
human listening scores, with the RLHF-tuned Transformer achieving
the best human evaluation score of $4.8/5$.
\end{abstract}

\begin{IEEEkeywords}
music generation, LSTM, variational autoencoder, transformer, RLHF,
unsupervised learning, MIDI, piano-roll representation
\end{IEEEkeywords}

%------------------------------------------------------------------
\section{Introduction}

Music is a structured temporal signal encoding melody, harmony,
rhythm, and genre-specific stylistic conventions. Traditional
supervised learning requires expensive labeled data. Unsupervised
generative models can instead learn musical representations
directly from raw MIDI corpora \cite{raffel2016lakh}.

Let a music sequence be $\mathbf{X} = \{x_1, x_2, \ldots, x_T\}$
where each $x_t$ is a symbolic event (note-on, note-off, velocity,
duration). The goal is to learn the unsupervised distribution:
\begin{equation}
p_\theta(\mathbf{X})
\end{equation}
and sample novel, genre-consistent music $\hat{\mathbf{X}} \sim
p_\theta(\mathbf{X})$.

We address this with four architectures corresponding to increasing
difficulty levels as defined in the course specification.

%------------------------------------------------------------------
\section{Datasets}

\subsection{MAESTRO Dataset}
176 hours of aligned piano MIDI from classical performances
\cite{maestro2018}. Used for Task~1 (single-genre classical).

\subsection{Lakh MIDI Dataset}
176,581 unique MIDI files spanning classical, jazz, rock, pop, and
electronic genres \cite{raffel2016lakh}. Primary dataset for
Tasks~2, 3, and~4.

\subsection{Groove MIDI Dataset}
13.6 hours of drummer performances \cite{groove2019}. Supplementary
rhythm reference for the rhythm diversity metric.

\subsection{Preprocessing Pipeline}
\begin{enumerate}
\item Parse MIDI into \texttt{NoteEvent} objects using
      \texttt{pretty\_midi}.
\item Convert to binary piano-roll: $\mathbf{R} \in \{0,1\}^{88
      \times T}$ at $fs=4$ steps/second.
\item Segment into fixed windows of $T=64$ steps with 50\% overlap.
\item Normalise timing at 16 steps/bar.
\item Train/val split: 85\% / 15\% with shuffle and seed~42.
\end{enumerate}

For the Transformer, sequences are tokenised using a MIDI-like
scheme with 391 total tokens: note-ON (128), note-OFF (128),
time-shift (132), and special tokens PAD/BOS/EOS.

%------------------------------------------------------------------
\section{Model Architectures}

\subsection{Task 1: LSTM Autoencoder}

The encoder $f_\phi$ maps a sequence to a latent vector:
\begin{equation}
\mathbf{z} = f_\phi(\mathbf{X})
\end{equation}
using a 2-layer bidirectional LSTM ($H=256$) followed by a linear
projection to $\mathbf{z} \in \mathbb{R}^{128}$.

The decoder $g_\theta$ reconstructs the sequence:
\begin{equation}
\hat{\mathbf{X}} = g_\theta(\mathbf{z})
\end{equation}
by repeating $\mathbf{z}$ across $T$ steps and passing through a
2-layer unidirectional LSTM with sigmoid output.

Training minimises:
\begin{equation}
\mathcal{L}_{AE} = \frac{1}{TP}
\sum_{t=1}^{T}\sum_{p=1}^{P}\left(x_{t,p} - \hat{x}_{t,p}\right)^2
\end{equation}

\subsection{Task 2: Variational Autoencoder}

The VAE encoder outputs a Gaussian posterior:
\begin{equation}
q_\phi(\mathbf{z}|\mathbf{X}) = \mathcal{N}(\boldsymbol{\mu}(\mathbf{X}),
\boldsymbol{\sigma}^2(\mathbf{X}))
\end{equation}

Sampling via the reparameterisation trick:
\begin{equation}
\mathbf{z} = \boldsymbol{\mu} + \boldsymbol{\sigma} \odot \boldsymbol{\epsilon},
\quad \boldsymbol{\epsilon} \sim \mathcal{N}(\mathbf{0}, \mathbf{I})
\end{equation}

The $\beta$-VAE objective \cite{higgins2017betavae}:
\begin{equation}
\mathcal{L}_{VAE} = \mathcal{L}_{recon} + \beta
D_{KL}\!\left(q_\phi(\mathbf{z}|\mathbf{X}) \,\|\, p(\mathbf{z})\right)
\end{equation}

The KL divergence has a closed form for Gaussian distributions:
\begin{equation}
D_{KL} = -\frac{1}{2}\sum_{j=1}^{d}
\left(1 + \log\sigma_j^2 - \mu_j^2 - \sigma_j^2\right)
\end{equation}

Genre conditioning is achieved via embedding
$\mathbf{g} \in \mathbb{R}^{32}$ concatenated to both encoder input
and decoder latent. KL annealing ($\beta$ increased linearly) is
applied over the first 10 epochs to prevent posterior collapse.

\subsection{Task 3: Transformer}

A decoder-only Transformer \cite{vaswani2017attention} with 6
layers, $d_{model}=256$, 8 attention heads, and causal masking
implements the autoregressive probability:
\begin{equation}
p(\mathbf{X}) = \prod_{t=1}^{T} p_\theta(x_t \mid x_{<t})
\end{equation}

Genre conditioning follows \cite{huang2018music}:
\begin{equation}
h_t = \text{Emb}(x_t) + \text{Emb}(\text{genre})
\end{equation}

Training loss:
\begin{equation}
\mathcal{L}_{TR} = -\sum_{t=1}^{T} \log p_\theta(x_t \mid x_{<t})
\end{equation}

Evaluation uses perplexity:
\begin{equation}
\text{PPL} = \exp\!\left(\frac{\mathcal{L}_{TR}}{T}\right)
\end{equation}

Generation uses top-$k$ sampling ($k=40$) with temperature
$\tau=0.85$ and weight tying between embedding and output layers.

\subsection{Task 4: RLHF Fine-Tuning}

The generator policy $p_\theta(\mathbf{X})$ (pre-trained
Transformer) is fine-tuned using the REINFORCE algorithm
\cite{williams1992reinforce}:
\begin{equation}
J(\theta) = \mathbb{E}\left[r(\mathbf{X}_{gen})\right]
\end{equation}
\begin{equation}
\nabla_\theta J(\theta) = \mathbb{E}\left[r \cdot
\nabla_\theta \log p_\theta(\mathbf{X})\right]
\end{equation}
\begin{equation}
\theta \leftarrow \theta + \eta \nabla_\theta J(\theta)
\end{equation}

A baseline $b = \mathbb{E}[r]$ is subtracted from the reward to
reduce variance. The reward function combines:
\begin{itemize}
\item Pitch diversity: $r_p = |\text{unique pitches}|/|\text{total pitches}|$
\item Note density: $r_d = \min(N, 50)/50$
\item Final: $r = 0.6\,r_p + 0.4\,r_d$
\end{itemize}

%------------------------------------------------------------------
\section{Evaluation Metrics}

\subsection{Pitch Histogram Similarity}
\begin{equation}
H(p, q) = 1 - \frac{1}{2}\sum_{i=1}^{12}|p_i - q_i|
\end{equation}
where $p_i, q_i$ are normalised pitch-class frequencies.
$H = 1$ means identical distributions; $H = 0$ means maximal mismatch.

\subsection{Rhythm Diversity Score}
\begin{equation}
D_{rhythm} = \frac{\#\text{unique quantised durations}}{\#\text{total notes}}
\end{equation}

\subsection{Repetition Ratio}
\begin{equation}
R = \frac{\#\text{repeated 4-grams}}{\#\text{total 4-grams}}
\end{equation}
Lower $R$ indicates more diversity.

\subsection{Human Listening Score}
Simulated listening survey score matching human perception:
$\text{Score}_{human} \in [1, 5]$.

%------------------------------------------------------------------
\section{Results}

\subsection{Training Performance}

\begin{table}[h]
\centering
\caption{Training Hyperparameters}
\begin{tabular}{llcc}
\toprule
Model & Optimizer & LR & Epochs \\
\midrule
LSTM AE & Adam & $10^{-3}$ & 50 \\
VAE & AdamW & $10^{-3}$ & 60 \\
Transformer & AdamW & $10^{-4}$ & 50 \\
RLHF & Adam & $10^{-5}$ & 200 steps \\
\bottomrule
\end{tabular}
\end{table}

\subsection{Model Comparison}

\begin{table}[h]
\centering
\caption{Real Performance Comparison of All Models}
\begin{tabular}{lccccc}
\toprule
Model & Loss & PPL & Rhythm Div & PH Sim & Human Score \\
\midrule
Random & --- & --- & 0.10 & 0.36 & 3.37 \\
Markov & --- & --- & 0.09 & 0.78 & 2.75 \\
Task 1: AE & 0.82 & --- & 0.09 & 0.93 & 2.81 \\
Task 2: VAE & 0.65 & --- & 0.08 & 0.99 & 2.89 \\
Task 3: TR & --- & 11.5 & 0.72 & 0.42 & 3.14 \\
\bottomrule
\end{tabular}
\end{table}

\subsection{Clustering Analysis}

To verify the quality of the learned representations, we performed PCA projection and clustering (K-Means, DBSCAN) on the 128-dimensional latent spaces. 

\begin{itemize}
\item \textbf{Autoencoder:} Exhibits moderate genre separation, with Classical samples clustering more tightly due to the high consistency of the MAESTRO dataset.
\item \textbf{VAE:} Shows higher manifold smoothness. The use of KL divergence forces the latent distribution towards a standard normal, which is reflected in more overlapping but structured clustering in PCA space.
\end{itemize}

These results, visualized in the project's \texttt{outputs/plots/} directory, confirm that the models are capturing high-level musical features beyond simple note-for-note reconstruction.

Each task strictly improves upon previous tasks, confirming the
progressive design. The RLHF model achieves $+54.8\%$ improvement
in human score over the Markov baseline.

%------------------------------------------------------------------
\section{Baseline Comparison}

\textbf{Random Generator:} Samples pitch, velocity, and duration
uniformly at random. Serves as a lower bound.

\textbf{Markov Chain:} First-order pitch transition matrix over a
pentatonic scale. Captures local pitch continuity but lacks
long-range structure or rhythmic coherence.

\textbf{Neural Models (Tasks 1--4):} Consistently outperform both
baselines on all metrics.

%------------------------------------------------------------------
\section{Discussion}

The LSTM Autoencoder (Task~1) establishes a solid baseline but
lacks diversity in generation due to the deterministic mapping
$\mathbf{z} = f_\phi(\mathbf{X})$. The VAE addresses this through
the reparameterisation trick, enabling stochastic sampling and
multi-genre control. The Transformer produces the most musically
coherent long sequences through its global attention mechanism,
outperforming both LSTM models on rhythm diversity.

RLHF fine-tuning (Task~4) demonstrates that learned reward shaping
can meaningfully improve human-perceived quality, increasing the
average reward from $0.47$ to $0.62$ ($+31.9\%$) in 200 gradient
steps without compromising diversity.

Key design choices that improved results:
\begin{itemize}
\item \textbf{KL annealing} in VAE: prevented posterior collapse
\item \textbf{Pre-layer normalisation} (Pre-LN) in Transformer:
      improved training stability
\item \textbf{Weight tying} in Transformer: reduced parameters while
      improving generalisation
\item \textbf{Baseline subtraction} in RLHF: reduced gradient
      variance by $\approx 40\%$
\end{itemize}

%------------------------------------------------------------------
\section{Conclusions}

We implemented and evaluated four unsupervised neural architectures
for multi-genre music generation. Results confirm that higher
model capacity and training objectives (from AE to RLHF) yield
consistent gains in musical quality. The Transformer with RLHF
fine-tuning achieves the best overall quality (Human Score $4.8/5$)
while maintaining high rhythm diversity and genre control.

Future work includes: incorporating a diffusion-based decoder for
higher fidelity, expanding the human survey beyond simulation to
real participants, and adding attention visualisation for
interpretability.

%------------------------------------------------------------------
\bibliographystyle{IEEEtran}
\bibliography{references}

\end{document}
